\chapter{Introduction generale}

\section{Contexte}

Le besoin metier vise une plateforme unifiee pour une academie de football, couvrant:
\begin{itemize}[leftmargin=*]
  \item la gestion administrative et operationnelle (joueurs, entraineurs, activites, paiements, notifications);
  \item l'experience utilisateur mobile (consultation, chat, profil, statistiques, recherche);
  \item une interface web admin dediee;
  \item un assistant conversationnel pour les questions frequentes des utilisateurs.
\end{itemize}

Cette combinaison impose une architecture multi-projets avec interconnexion claire, separation des responsabilites et securisation des acces.

\section{Problematique}

Les principales questions techniques traitees dans ce rapport sont:
\begin{itemize}[leftmargin=*]
  \item comment organiser le backend pour supporter a la fois API REST, temps reel et controle d'acces par roles;
  \item comment structurer un frontend Flutter qui sert mobile et web admin sans duplication inutile;
  \item comment integrer un chatbot local independant et interopere avec l'application mobile;
  \item comment preparer un deploiement robuste et evolutif.
\end{itemize}

\section{Demarche}

La demarche retenue est basee sur l'analyse documentaire et du code existant:
\begin{itemize}[leftmargin=*]
  \item extraction des technologies, configurations et flux depuis les fichiers de configuration et README;
  \item cartographie des routes, couches logicielles et dependances;
  \item consolidation en chapitres exploitables pour un rapport academique.
\end{itemize}

\section{Structure du rapport}

Le rapport est organise en chapitres progressifs:
\begin{enumerate}[leftmargin=*]
  \item sources documentaires exploitees;
  \item architecture globale de la solution;
  \item etude du backend;
  \item etude du frontend;
  \item etude du chatbot;
  \item integration, deploiement, qualite et perspectives.
\end{enumerate}
